\documentclass{scrartcl}

\usepackage[T1]{fontenc}
\usepackage[utf8]{inputenc}
\usepackage{times}
\usepackage{geometry}
\usepackage{lastpage}
\usepackage{fancyhdr}
\usepackage{hyperref}
\usepackage{enumitem}
\usepackage{microtype}

% Version number
\def\version{0.1}

\geometry{letterpaper,margin=1in}
\setlength{\parindent}{0pt}
\setlength{\parskip}{0.6em}
\setlist{nosep,leftmargin=*}
\hypersetup{hidelinks,pdfauthor={CSU, Chico Computer Science Department},pdftitle={Department Constitution and Bylaws}}

\title{California State University, Chico\\Computer Science Department\\Constitution and Bylaws}
\subtitle{Version \version}
\date{}

\pagestyle{fancy}
\fancyhf{}
\fancyhead[L]{\fontsize{8}{10}\selectfont CSU, Chico Computer Science Department Constitution and Bylaws}
\fancyhead[R]{\fontsize{8}{10}\selectfont Version \version\\Page \thepage\ of \pageref{LastPage}}

\begin{document}

\begin{titlepage}
\clearpage
\maketitle
\thispagestyle{empty}
\end{titlepage}

\tableofcontents
\clearpage

\section*{SECTION I PURPOSE (Sam)}
\phantomsection
\addcontentsline{toc}{section}{SECTION I PURPOSE (Sam)}


The Department of Computer Science Constitution and Bylaws have been developed in accordance with the University’s Personnel Policies and Procedures (FPPP) and the Collective Bargaining Agreement (CBA). The section outlines standards, policies, and procedures that consider the unique qualities and needs of the Department of Computer Science as a professional program within the University.

All faculty members must contribute to developing the department and its programs. The tasks and duties around these responsibilities will be given significant weight in the retention, tenure and promotion process.

\section*{SECTION II DEPARTMENT MEMBERSHIP}
\phantomsection
\addcontentsline{toc}{section}{SECTION II DEPARTMENT MEMBERSHIP}


Membership shall be awarded in the following categories:

\begin{enumerate}
\item All tenured or tenure-track professors (assistant professors, associate professors, and full professors) holding at least half-time appointments in the Department.
\item Lecturer membership shall be awarded to all adjunct faculty and lecturers holding appointments in the Department, including retired lecturers with concurrent teaching assignments.
\item Retired faculty membership shall be granted to all Emeritus faculty and those within the Faculty Early Retirement Program (FERP).
\item Staff membership shall be granted to all non-teaching staff holding appointments of 10/12 or greater in the Department.
\end{enumerate}

Membership in the Department of Computer Science shall not lapse because of sabbatical leave or any other approved leave of absence.

\section*{SECTION III STANDING, AD HOC, AND OTHER DEPARTMENTAL COMMITTEES}
\phantomsection
\addcontentsline{toc}{section}{SECTION III STANDING, AD HOC, AND OTHER DEPARTMENTAL COMMITTEES}


\textbf{Section 1.} Committees of the whole are preferred. When a committee of the whole is deemed appropriate by the Department of Computer Science members and a smaller committee is formed, all interested and eligible members should be encouraged to stand for election. The election of committee members shall be by a simple majority vote.

\textbf{Section 2.} Each committee shall make its recommendations and reports to the Department of Computer Science unless specifically directed to report to the Department Chair.

\textbf{Section 3.} Each committee shall meet at the call of its chair if one is elected when it is formed. The Department Chair will call the first meeting if no chair has been selected. A chair will be elected at the first meeting if one does not already exist.

\textbf{Section 4.} Each committee will conduct its business according to the rules established by the Department when the committee is formed. If the Department establishes no such rules, the committee will conduct its business according to the rules agreed to by the committee members.

\textbf{Section 5.} Standing committees such as CSCI Curriculum, CINS Curriculum, Graduate Curriculum, Assessment, etc., may only be formed by a two-thirds vote of the Department.

\textbf{Section 6.} Ad Hoc Committees such as Search/Hiring, Textbook, and Student Recruitment may be formed by a simple majority vote.

\textbf{Section 7.} The chair of a standing committee, even when it is a committee of the whole, need not be the Department Chair.

\textbf{Section 8.} The term for each committee member will be one year unless it is decided otherwise by majority vote of the Department.

\textbf{Section 9.} Members of committees may be recalled by a two-thirds vote of the Department.

\section*{SECTION IV HIRING}
\phantomsection
\addcontentsline{toc}{section}{SECTION IV HIRING}


\subsection*{A. Tenure Track (Probationary) Faculty}
\phantomsection
\addcontentsline{toc}{subsection}{A. Tenure Track (Probationary) Faculty}


\begin{enumerate}
\item Equivalency Qualifications (Provost approval required)
\item Procedures
\item Service Credit
\end{enumerate}

\subsection*{B. Temporary Faculty}
\phantomsection
\addcontentsline{toc}{subsection}{B. Temporary Faculty}


\begin{enumerate}
\item Rationale for Temporary Appointments
\item Appointment Procedures
\item University Appointment Standards for Lecturer Ranges
\end{enumerate}

\section*{SECTION V MEETINGS}
\phantomsection
\addcontentsline{toc}{section}{SECTION V MEETINGS}


\begin{enumerate}
\item Departmental Meetings

\begin{enumerate}
\item Departmental meetings will be held weekly during the fall and spring semesters, with the exception of finals week.
\item Departmental meetings during the fall and spring semesters will be held on Wednesdays at 2:00 pm.
\begin{enumerate}[label=\alph*.]
\item The department chair will ensure no faculty are scheduled during this time.
\end{enumerate}
\item As scheduled by the department chair, a departmental meeting will be held the week before the start of the fall and spring semesters. Typically, this meeting will occur on the same day as the College Semester Kick-off Meeting.
\item The department's administrative support coordinator (ASC) will take minutes for every departmental meeting.
\begin{enumerate}[label=\alph*.]
\item In the absence of the ASC, minutes will be recorded by another department member in attendance appointed by the department chair.
\end{enumerate}
\item The department’s ASC will maintain minutes for 5 full academic years.
\item Additional or special departmental meetings can be scheduled as needed.
\item The department chair sets the agenda for every department meeting.
\begin{enumerate}[label=\alph*.]
\item Committee chairs will be given the opportunity to introduce items onto the agenda pertaining to their specific committee
\item All department members will be given the opportunity to introduce items onto the agenda that pertain to the department as a whole and are not under the purview of a standing committee.

\end{enumerate}
\end{enumerate}
\item Committee Meetings

\begin{enumerate}
\item Committee meetings will be held as needed as determined by the committee chair.
\item Committee meetings will meet only during the fall and spring semesters.
\item Committee meetings will be held on Mondays at 2:00 pm.
\begin{enumerate}[label=\alph*.]
\item The department chair will ensure no faculty are scheduled during this time.
\end{enumerate}
\item The chair or appointed committee member will take minutes for every meeting.
\item The committee chair will maintain minutes for 3 full academic years.
\end{enumerate}
\end{enumerate}

\section*{SECTION VI VOTING}
\phantomsection
\addcontentsline{toc}{section}{SECTION VI VOTING}


\paragraph{A. Tenure, Tenure Track, and FERP Faculty}


\begin{enumerate}
\item All tenured and tenure-track faculty have the right to vote.
\item Personnel with a right to vote have the right to a proxy.

\end{enumerate}
\paragraph{B. Lecturer Faculty}


\begin{enumerate}
\item All lecturer faculty with a current teaching assignment can vote for Department Chair elections.
\item Lecturer faculty who hold a teaching assignment of 6 WTUS or more have the right to a vote in all other matters not outlined in Section VI.B.1

\end{enumerate}
\paragraph{C. Staff Members}


\begin{enumerate}
\item Non-teaching staff members holding appointments of 10/12 or greater in the department have the right to vote for Department Chair elections.
\item Non-teaching staff members, regardless of appointment length, do not have the right to vote, with the exception of rights outlined in Section VI.C.1.
\item Staff members do not have the right to a proxy.

\end{enumerate}
\paragraph{D. Voting Rights While on Leave}


\begin{enumerate}
\item As Section VI outlines, any department member with the right to vote retains that right when on paid leave.
\item Any department member with the right to a vote as outlined by Section VI loses their right to a vote when on unpaid leave.
\item As Section VI outlines, any department member has the right to a proxy. However, they lose this right when on leave.

\end{enumerate}
\paragraph{E. Voting Proxies}


\begin{enumerate}
\item A faculty member may only be given a proxy for at most one individual with a right to vote as defined in Section VI.

\end{enumerate}
\paragraph{F. Quorum}


\begin{enumerate}
\item A simple majority of the voting membership is defined in Section VI.
\end{enumerate}

\section*{SECTION VII NOMINATION AND DUTIES OF OFFICERS (Bryan)}
\phantomsection
\addcontentsline{toc}{section}{SECTION VII NOMINATION AND DUTIES OF OFFICERS (Bryan)}


Chair, vice chair, graduate coordinator, lab coordinator

\subsection*{Department Chair}
\phantomsection
\addcontentsline{toc}{subsection}{Department Chair}


\textbf{Section 1.} The Department Chair is a member of the Department of Computer Science faculty elected to perform administrative functions and provide leadership in developing the strongest possible academic program. The duties of a department chair should be considered a full-time role.

\textbf{Section 2.} The minimum responsibilities of the Department Chair are detailed in FPPP.

\textbf{Section 3.} Additional responsibilities of the Department Chair include the following:
\begin{itemize}
\item Establish and maintain the department’s Advisory Board and its schedule of meetings
\item Schedule and run department meetings
\item Chair the department RTP committee or serve as chair-level review
\item Overall responsibility for class schedule
\item Overall responsibility for teaching assignments
\item Overall responsibility for program accreditation
\item Overall responsibility for advising and course substitutions
\item Hiring of lectures
\item Justification for tenure track positions
\item Manage program assessment, surveys, and student learning outcomes
\item Manage Choose Chico Day
\item Manage Chico Preview Day
\item Manage Alumni in the Classroom Day
\item Direct content for the department website
\item Oversee Internships and documentation for technical elective credit
\item Oversee the department’s budget and spending
\item Oversee the department’s labs and ensure lab safety
\item Approve department equipment
\item Oversee program curricula and catalog information
\item Develop and oversee department marketing materials
\item Serve as the department representative on the ECC Curriculum Committee
\item Attend ECC Chair-Director meetings
\item Attend University Chairs Council meetings
\item Attend University Extended Leadership Forum (ELF) meetings
\item Attend University MPP/Chair meetings
\item Evaluation of department ASC
\end{itemize}

\textbf{Section 4.} The following procedures will be used for selection of the Department Chair.
\begin{enumerate}
\item Department Chair elections are to be carried out during the last academic year of the the current term of the incumbent department chair.
\item The Department faculty will elect a Chair nominating committee consisting of three tenured or tenure-track faculty from the Department.
\item The committee will review the Chair selection criteria (Article VI, Section 5).
\item The committee will present the Chair selection criteria to the Department. The The department will vote to accept or reject any proposed revisions.
\item The committee will call for nominations for the Department Chair. One week following the call for nominations, the committee will submit to the Department the names of all qualified candidates.
\item A forum will be held that provides the candidates an opportunity to present their candidacy to the Department with opportunities for questions and answers.
\item The department will then elect the Chair nominee by secret ballot. In the event of a single nominee, the Department may elect by acclamation. The term of office shall be three years unless otherwise stated before balloting. Ballot procedures will conform to College practices regarding secret ballots. The nominee must receive a majority of the votes cast.
\item The Chair nominating committee will forward the name of the elected candidate to the Dean of the College of Engineering, Computer Science \& Construction Management within ten days following this election.
\item The Chair nominating committee will ask the Dean to forward the faculty selection with recommendation to the Provost. Appointment of the Chair will be made by the Provost as chief instructional officer. The commencement date will be negotiated by the nominee with the Dean.
\end{enumerate}

\textbf{Section 5.} Department Chair Selection Criteria
\paragraph{A. Professional Qualifications}

\begin{enumerate}
\item The candidate should have considerable and varied experience in computer science, computer information systems, or related specialization.
\item It is desirable that the candidate hold appropriate professional credentials in his/her field.
\end{enumerate}
\paragraph{B. Academic Preparation}

\begin{enumerate}
\item It is highly desirable for the candidate to have at least one advanced degree with an earned doctoral degree preferred in computer science, computer information systems, or related specialization.
\item The candidate must have considerable teaching experience in an accredited program, preferably by ABET or another recognized accreditation body. A satisfactory record of classroom teaching and course development is required.
\item A record of significant research activities and recent publications in refereed technical journals would be considered an asset.
\item The candidate should demonstrate the ability to remain current in his/her area of specialization.

\end{enumerate}
\paragraph{C. Administrative Ability}

\begin{enumerate}
\item The candidate should have the capacity to make decisions and to pursue them to conclusion.
\item The candidate should have the ability to guide the curricular development, assessment, and accreditation of the computer science and computer information systems programs.
\item The candidate should be able to utilize to the best advantage whatever resources are at his/her disposal. The candidate should demonstrate the ability to obtain additional resources from government and/or industrial sources.
\item The candidate should have demonstrated administrative ability, managing and directing projects and/or personnel.
\item The candidate should be able to attract and maintain a highly qualified faculty.
\item The candidate should be capable of initiating adequate budget proposals and making effective use of funds allocated to the department.
\item The candidate should be capable of carrying out all functions implied by the Responsibilities of Department Chairs in the FPPP.

\end{enumerate}
\paragraph{D. General}

\begin{enumerate}
\item Candidates will be considered on an overall basis. A deficiency in one area may be offset by superior qualifications in other areas.
\item The numbering of items above is not intended to imply priority order.
\end{enumerate}

\textbf{Section 6.} The Department Chair term will be a three-year appointment.

\textbf{Section 7.} The Department Chair may be recalled by a two-thirds vote.

\subsection*{The Department Vice-Chair}
\phantomsection
\addcontentsline{toc}{subsection}{The Department Vice-Chair}


\textbf{Section 1.} At the Department Chair's discretion, the Department Chair's position may be split to include the position of Department Vice-Chair. This would result in a quarter-time Vice-Chair and three-quarter Department Chair unless a larger split is agreed upon.

\textbf{Section 2.} The Department Vice-Chair is a member of the faculty of the Department of Computer Science who is nominated by the Department Chair.

\textbf{Section 3.} The Vice-Chair nominee must satisfy the selection criteria described for the Department Chair in Article VI, Section 5. The Vice-Chair nominee must be approved by the Department vote.

\textbf{Section 4.} The Department Chair will forward the name of the elected Vice-Chair to the Dean of the College of Engineering, Computer Science \& Construction Management by the Department Chair within ten days following this election.

\textbf{Section 5.} The Department Chair will ask the Dean to forward the Vice-Chair selection with a recommendation to the Provost. Appointment of the Vice-Chair will be made by the Provost as chief instructional officer. The date of the commencement will be negotiated by the Vice-Chair with the Dean.

\textbf{Section 6.} The Vice-Chair and Department Chair must agree on a division of the responsibilities described in Article VI, Sections 1-3.

\textbf{Section 7.} The Vice-Chair term will be a one-year appointment, renewable at the Department Chair’s discretion.

\textbf{Section 8.} The Department Vice-Chair may be recalled by a two-thirds vote.

\subsection*{Graduate Coordinator and Online Graduate Coordinator}
\phantomsection
\addcontentsline{toc}{subsection}{Graduate Coordinator and Online Graduate Coordinator}


\textbf{Section 1.} At the Department Chair's discretion, the Department may appoint a graduate coordinator. This would be a half-time appointment. The Department Chair will appoint a graduate coordinator for in-person and online graduate programs. These are separate roles that oversee these two separate programs.

\textbf{Section 2.} The Department Graduate Coordinator is a member of the Department of Computer Science faculty nominated by the Department Chair.

\textbf{Section 3.} The Graduate Coordinator nominee must satisfy the selection criteria described for the Department Chair in Article VI, Section 5. The Graduate Coordinator nominee must be approved by the Department vote.

\textbf{Section 4.} The Department Chair will forward the name of the elected Graduate Coordinator to the Dean of the College of Engineering, Computer Science \& Construction Management by the Department Chair within ten days following this election.

\textbf{Section 5.} The Department Chair will ask the Dean to forward the Graduate Coordinator selection with a recommendation to the Provost. Appointment of the Graduate Coordinator will be made by the Provost as chief instructional officer. The date of the commencement will be negotiated by the Graduate Coordinator with the Dean.

\textbf{Section 6.} The Graduate Coordinator will be responsible for:

\begin{itemize}
\item Reviewing applications of students who have applied to the program
\item Determining which portion of students should be accepted in communication with the Department Chair of acceptance number to target.
\item Writing acceptance letters and prerequisite requirements incoming students will need to satisfy to be classified.
\item Advising current graduate students and approving their course programs.
\item The online graduate coordinator will schedule the courses offered and find faculty to teach them for every cohort or session.
\end{itemize}

\textbf{Section 7.} The Graduate Coordinator term will be a one-year appointment, renewable at the Department Chair’s discretion.

\textbf{Section 8.} The Department Graduate Coordinator may be recalled by a two-thirds vote.

\subsection*{Lab Coordinator}
\phantomsection
\addcontentsline{toc}{subsection}{Lab Coordinator}


\textbf{Section 1.} At the Department Chair's discretion, the Department may appoint a Lab Coordinator. This would be a quarter-time appointment.

\textbf{Section 2.} The Department Lab Coordinator is a member of the Department of Computer Science faculty nominated by the Department Chair.

\textbf{Section 3.} The Lab Coordinator nominee must satisfy the selection criteria described for the Department Chair in Article VI, Section 5. The Department must approve the Lab Coordinator nominee by vote.

\textbf{Section 4.} The Department Chair will forward the name of the elected Lab Coordinator to the Dean of the College of Engineering, Computer Science \& Construction Management by the Department Chair within ten days following this election.

\textbf{Section 5.} The Department Chair will ask the Dean to forward the Lab Coordinator selection with a recommendation to the Provost. The Provost, as chief instructional officer, will appoint the Lab Coordinator. The Lab Coordinator will negotiate the commencement date with the Dean.

\textbf{Section 6.} The Lab Coordinator will be responsible for:

\begin{itemize}
\item Interfacing with ESYS and ITSS regarding department needs involving:
\begin{itemize}
\item Lab machines
\begin{itemize}
\item Packages and dependencies needed
\end{itemize}
\item Hosted resources
\begin{itemize}
\item Servers hosted in the collab space in Butte
\item Servers hosted by ESYS in the cloud
\item Student shared storage
\item Remote Linux server
\end{itemize}
\item VPN support documentation for Linux
\item Submitting ITPR tickets or following up on tickets submitted by other faculty
\end{itemize}
\item Managing department automated exercises assessment platform
\begin{itemize}
\item Installing updates
\item Handling security vulnerability notifications
\item Handling bugs identified by faculty using the program
\item Helping get faculty setup and using the system
\end{itemize}
\item Managing department server resources
\begin{itemize}
\item Managing access to virtual server space for faculty
\item Managing access to a coprocessor server, if the department has one, for faculty and students in courses needing it.
\end{itemize}
\item Serve as the chair of the department lab committee
\end{itemize}

\textbf{Section 7.} The Lab Coordinator term will be a one-year appointment, renewable at the Department Chair’s discretion.

\textbf{Section 8.} The Department Lab Coordinator may be recalled by a two-thirds vote.

\section*{SECTION VIII ORDER OF MEETING BUSINESS: (Sam)}
\phantomsection
\addcontentsline{toc}{section}{SECTION VIII ORDER OF MEETING BUSINESS: (Sam)}


The Department Chair shall call meetings as the Chair deems necessary. The Department of Computer Science membership may also request meetings. Upon receiving such a request from at least two members, the Department Chair shall call a meeting within ten days.

The Department Chair will make every effort to ensure that meeting times do not conflict with members' academic responsibilities.

All members of the Department of Computer Science in residence will be notified of the time and agenda at least two days prior to every meeting. Minutes shall be made available to members in a timely fashion.

Meetings are preferred to be run informally, consistent with a tradition of collegiality. However, if a conflict over procedural matters persists, Robert's Rules of Order shall be the competent authority.

Meetings where decisions are made by vote (using Robert’s Rules as necessary) require a quorum and a clear agenda item or prior motion to discuss. A departmental quorum requires a majority of members to be present, including the chair, assistant chair, and any committee chairs bringing a matter to the department for a vote.  Members may delegate their vote if they cannot be present, but must notify the chair in advance for the delegated vote to count.

\section*{SECTION IX TRAVEL FUNDS (Sam)}
\phantomsection
\addcontentsline{toc}{section}{SECTION IX TRAVEL FUNDS (Sam)}


Travel funds from the department professional development account may be requested by any department member using a simple Google form or similar for the purposes of professional development including:

\begin{enumerate}
\item Conference attendance, with priority for attendance where a paper is presented by the applicant or a student mentored by the applicant
\item Training and certification courses
\item Course and curriculum development seminars
\item Industry advisory council activities
\item Competition, club and society events if the sponsoring organization lacks funding (yes/no/maybe)
\end{enumerate}

\section*{SECTION X COURSE SCHEDULING: (Bryan)}
\phantomsection
\addcontentsline{toc}{section}{SECTION X COURSE SCHEDULING: (Bryan)}


\paragraph{A. Academic year}

\paragraph{B. Winter Intersession}

\paragraph{C. Summer Session(s)}


\section*{SECTION XI FACULTY LEAVES (Bryan)}
\phantomsection
\addcontentsline{toc}{section}{SECTION XI FACULTY LEAVES (Bryan)}


\section*{SECTION XII GRADUATE PROGRAM CONSIDERATIONS: (Sam)}
\phantomsection
\addcontentsline{toc}{section}{SECTION XII GRADUATE PROGRAM CONSIDERATIONS: (Sam)}


Department compliance with graduate studies

Graduate studies requires all graduate programs to have a “culminating” activity in addition to required credit hours and courses for students to become graduation candidates and in the Computer Science department, this can be satisfied by:

\begin{enumerate}
\item CSCI 693 - Research Methods
\item CSCI 699P - Research Project
\item CSCI 699T - Thesis
\end{enumerate}

The CSCI 693 option is offered to all students as a class taught by faculty members based upon interest and assignment by the department chair.  The CSCI 699P and 699T options are based upon mutual interest between a student with a project or thesis proposal and a potential advisor/mentor. Advisors for the project or thesis option will receive 0.5 WTU for their efforts over one or two semesters based upon a proposed schedule that must be submitted no later than 1 month prior to the end of the semester prior to the start of this work.

\paragraph{A. Thesis Exams}


Graduate courses including thesis and project options must include assessment in the form of exams that can be written, oral presentations, or submitted documents and code. Generally a thesis exam is oral (by committee) and code and documentation are submitted with the written thesis to graduate studies. The thesis and project exams must meet all graduate studies requirements.

All exams for all graduate courses of study should be administered by a faculty member or department member proctor and students are expected to follow all California State University academic honesty policies published for the system and the Chico campus as well as any specific course policies the instructor may have included in their syllabus.

Written project and thesis work will be assessed for plagiarism using university sanctioned tools such as Turn-it-in or similar supplied by the university and all work is expected to be the original work of the thesis or project student or properly cited as required by graduate studies.  Code is likewise must be original work, but may include libraries, generated code, and code re-used as allowed by the student’s advisor, but also must be properly cited and adhere to licensing requirements.

\paragraph{B. Theses and Projects}


For CSCI 699T students electing to complete a thesis instead of CSCI 693 Research Methods or CSCI 699P Project, must submit a detailed thesis proposal during the semester before they start work for credit at least one month before the end of that semester and so that all department faculty members have ample time to review the proposal (at least two full weeks). During this time, all faculty members can raise concerns, objections, or make suggestions for improvement that require the student and advisor to respond.  The proposal, reviews, and rebuttals should be accessible to all faculty members.

Thesis proposals must include an advisor and committee members from the Computer Science faculty members (as required by graduate studies) identified at the time of proposal submission. In extenuating circumstances (e.g., committee member illness) a committee member substitution can be made with approval of the current committee. All theses should be completed within two semesters of the accepted proposal, but may straddle a summer break, and in extenuating circumstances such as student illness or resource issues can be extended only if the thesis committee agrees to the extension unanimously.

If a graduate student does not complete their thesis including thesis oral exam and successful submission of their written thesis to graduate studies, they can take up to one more semester to complete their work with unanimous approval by all thesis committee members.

If a graduate student electing to complete a thesis does not pass their oral exam, is unable to complete the written thesis, or agrees with their advisor to drop the thesis enrollment and objective, they may still complete the graduate program in the semester following by enrolling in CSCI 693 Research Methods.

\paragraph{C. Faculty Committee Composition}


Thesis and project committees for CSCI 699P and 699T shall meet the minimum requirements of graduate studies or exceed and shall have a Computer Science department member as advisor with at least one other Computer Science department member (other members may come from outside the department with unanimous approval by the other committee members).

\section*{SECTION XIII NEW FACULTY MENTORSHIP}
\phantomsection
\addcontentsline{toc}{section}{SECTION XIII NEW FACULTY MENTORSHIP}


\paragraph{A. Senior faculty classroom observations by new faculty}

\textbf{Purpose:} The purpose of this document is to formalize the expectations of new faculty members in respect to observing classes within the department. New faculty members may benefit personally from such observations in the following ways: (1) be exposed to different methods of instruction, (2) explore the different course offerings of the department, and (3) strengthen working relationships with colleagues in the department.

\textbf{Proposed Guidelines:} Any new faculty member who is hired to the computer science department, beginning in Fall 2025, will be expected to observe two classes in the department before their first Retention, Tenure, and Promotion (RTP) performance review.

They must visit two different courses taught by two different faculty members who have been with the department for longer than the newly hired faculty member. The newly hired faculty member should meet with the course instructor to debrief about their observations following the class visit. The discussion should touch on the pedagogical practices observed and how the new faculty member is adjusting to joining the department.

The non-new faculty member should reach out to the new faculty member with times that would be good for a class visit.

There will not be a required write-up for the observations. However, new faculty members are highly encouraged to incorporate what they learned from the experience into their Personnel Action File during the RTP process.

Failure to complete these class observations will render a faculty member ineligible for "meeting  expectations" in the instructional effectiveness area of the RTP performance review, unless the department agrees to waive this expectation for a particular faculty member based on extraneous circumstances. By the 2nd year evaluation, faculty are expected to have evidence (such as email correspondence) of their scheduled class visits in order to meet expectations. There are no requirements for the 4th or 6th year evaluations, but additional visits will allow the faculty member to achieve “exceeding expectations” within the respective RTP category.

\paragraph{B. Faculty Mentorship Recommendations}


\section*{SECTION XIV ETHICS AND STANDARDS}
\phantomsection
\addcontentsline{toc}{section}{SECTION XIV ETHICS AND STANDARDS}


All Department of Computer Science members must comply with the conflict of interest section as described in FPPP.

All members of the Department of Computer Science must observe the Code of Ethics as described in FPPP Appendix II as well as the Executive Memorandum 09-008 regarding Nepotism.

Amendments to this Constitution may be proposed by any member of the Department of Computer Science. The proposal must be seconded by another member and made available to the Department at least ten days before the vote.

Amendments to this Constitution shall require approval by a two-thirds vote.

\section*{SECTION XV AMENDMENTS (Bryan)}
\phantomsection
\addcontentsline{toc}{section}{SECTION XV AMENDMENTS (Bryan)}


\textbf{Section 1.} Amendments to this Constitution may be proposed by any member of the Department of Computer Science. Another member must second the proposal, and the proposal must be made available to the Department at least ten days before the vote.

\textbf{Section 2.} Amendments to this Constitution shall require approval by a two-thirds vote.

\end{document}
